% paper23-evidence-routing.tex — Mixed Evidence Routing: Theory-Aware Dispatch
% of Verification Conditions
% Paper 23 of the Judgment Geometry series.
% Compile: pdflatex paper23-evidence-routing && pdflatex paper23-evidence-routing
\documentclass[11pt]{article}
\usepackage{lmodern}
% jugeo-common.tex — shared preamble for all Judgment Geometry papers
% Include via: % jugeo-common.tex — shared preamble for all Judgment Geometry papers
% Include via: % jugeo-common.tex — shared preamble for all Judgment Geometry papers
% Include via: \input{jugeo-common}

% ── Packages ────────────────────────────────────────────────────────────
\usepackage{amsmath,amssymb,amsthm,mathtools}
\usepackage{tikz-cd}
\usepackage{listings}
\usepackage{xcolor}
\usepackage{xspace}
\usepackage{booktabs}
\usepackage{multirow}
\usepackage{enumitem}
\usepackage[expansion=false]{microtype}
\usepackage{hyperref}
\usepackage{cleveref}
% \usepackage{thmtools}  % disabled: not available in this TeX installation
\usepackage{float}
\usepackage{caption}
\usepackage{subcaption}
\usepackage{tabularx}
\usepackage{stmaryrd}       % \llbracket, \rrbracket
\usepackage{bbm}            % \mathbbm{1}

% ── Page geometry ───────────────────────────────────────────────────────
\usepackage[margin=1in]{geometry}

% ── Theorem environments ────────────────────────────────────────────────
\theoremstyle{plain}
\newtheorem{theorem}{Theorem}[section]
\newtheorem{lemma}[theorem]{Lemma}
\newtheorem{proposition}[theorem]{Proposition}
\newtheorem{corollary}[theorem]{Corollary}

\theoremstyle{definition}
\newtheorem{definition}[theorem]{Definition}
\newtheorem{example}[theorem]{Example}
\newtheorem{construction}[theorem]{Construction}

\theoremstyle{remark}
\newtheorem{remark}[theorem]{Remark}
\newtheorem{notation}[theorem]{Notation}

% ── Python listings style ──────────────────────────────────────────────
\definecolor{jugeoBlue}{HTML}{2563EB}
\definecolor{jugeoGreen}{HTML}{059669}
\definecolor{jugeoRed}{HTML}{DC2626}
\definecolor{jugeoPurple}{HTML}{7C3AED}
\definecolor{jugeoGray}{HTML}{6B7280}
\definecolor{jugeoBg}{HTML}{F8FAFC}

\lstdefinestyle{jugeo-python}{
  language=Python,
  basicstyle=\ttfamily\small,
  keywordstyle=\color{jugeoBlue}\bfseries,
  stringstyle=\color{jugeoGreen},
  commentstyle=\color{jugeoGray}\itshape,
  numberstyle=\tiny\color{jugeoGray},
  backgroundcolor=\color{jugeoBg},
  numbers=left,
  numbersep=6pt,
  frame=single,
  framerule=0.4pt,
  rulecolor=\color{jugeoGray!40},
  breaklines=true,
  breakatwhitespace=true,
  showstringspaces=false,
  tabsize=4,
  xleftmargin=1.5em,
  framexleftmargin=1.5em,
  morekeywords={dataclass,frozen,field,Enum,IntEnum,auto,Any,
                Mapping,Sequence,Optional,match,case},
  literate={→}{{$\to$}}1 {←}{{$\leftarrow$}}1
           {≤}{{$\leq$}}1 {≥}{{$\geq$}}1 {≠}{{$\neq$}}1
           {∈}{{$\in$}}1 {∀}{{$\forall$}}1 {∃}{{$\exists$}}1
           {⊕}{{$\oplus$}}1 {⊖}{{$\ominus$}}1,
  escapeinside={(*@}{@*)},
}
\lstset{style=jugeo-python}

\lstdefinestyle{jugeo-output}{
  basicstyle=\ttfamily\small,
  backgroundcolor=\color{jugeoBg},
  frame=single,
  framerule=0.4pt,
  rulecolor=\color{jugeoGray!40},
  breaklines=true,
  showstringspaces=false,
  xleftmargin=1.5em,
  framexleftmargin=0.5em,
  numbers=none,
}

% ── JuGeo-specific macros ──────────────────────────────────────────────

% Judgment 8-tuple
\newcommand{\J}{\mathcal{J}}
\newcommand{\Jud}[8]{(#1,\,#2,\,#3,\,#4,\,#5,\,#6,\,#7,\,#8)}
\newcommand{\JudShort}[1]{\J_{#1}}

% Components
\newcommand{\coord}{\mathsf{c}}            % coordinate
\newcommand{\prop}{\varphi}                 % proposition
\newcommand{\carrier}{\mathcal{A}}          % carrier type
\newcommand{\evid}{\mathcal{E}}             % evidence bundle
\newcommand{\oblig}{\mathcal{O}}            % obligations
\newcommand{\obstr}{\mathcal{B}}            % obstructions
\newcommand{\trust}{\mathcal{T}}            % trust annotation
\newcommand{\prov}{\Pi}                     % provenance

% Trust algebra
\newcommand{\Talg}{\mathfrak{T}}            % trust ordered algebra
\newcommand{\Eadm}{\mathcal{E}_{\mathrm{adm}}}
\newcommand{\tleq}{\preceq}                 % trust ordering
\newcommand{\tjoin}{\oplus}                 % conservative join
\newcommand{\tatten}{\ominus}               % attenuation
\newcommand{\tpromote}{\uparrow_\pi}        % promotion
\newcommand{\tdemote}{\downarrow_\chi}       % demotion

% Trust tiers
\newcommand{\Tcontra}{\ensuremath{\mathsf{CONTRADICTED}}}
\newcommand{\Tunver}{\ensuremath{\mathsf{UNVERIFIED}}}
\newcommand{\Tcopilot}{\ensuremath{\mathsf{COPILOT}}}
\newcommand{\Toracle}{\ensuremath{\mathsf{ORACLE}}}
\newcommand{\Truntime}{\ensuremath{\mathsf{RUNTIME}}}
\newcommand{\Tsolver}{\ensuremath{\mathsf{SOLVER}}}
\newcommand{\Tproof}{\ensuremath{\mathsf{PROOF}}}

% Site and geometry
\newcommand{\Site}{\mathcal{S}}             % semantic site
\newcommand{\Cov}{\mathrm{Cov}}            % covering family
\newcommand{\Ob}{\mathrm{Ob}}              % objects
\newcommand{\Mor}{\mathrm{Mor}}            % morphisms
\newcommand{\res}{\mathrm{res}}            % restriction
\newcommand{\Cech}{\check{C}}              % Čech complex
\newcommand{\Hcoh}[1]{H^{#1}}             % cohomology group

% Descent
\newcommand{\descent}{\mathrm{desc}}
\newcommand{\glue}{\mathrm{glue}}
\newcommand{\overlap}[2]{#1 \cap #2}

% Semantic moves
\newcommand{\VERIFY}{\textsc{Verify}}
\newcommand{\CONSTRUCT}{\textsc{Construct}}
\newcommand{\NORMALIZE}{\textsc{Normalize}}
\newcommand{\GLUE}{\textsc{Glue}}
\newcommand{\RESTRICT}{\textsc{Restrict}}
\newcommand{\TRANSPORT}{\textsc{Transport}}
\newcommand{\REPAIR}{\textsc{Repair}}
\newcommand{\PROMOTE}{\textsc{Promote}}
\newcommand{\CHALLENGE}{\textsc{Challenge}}
\newcommand{\ESCALATE}{\textsc{Escalate}}

% Fragments
\newcommand{\QFLIA}{\texttt{QF\_LIA}}
\newcommand{\QFLRA}{\texttt{QF\_LRA}}
\newcommand{\QFBV}{\texttt{QF\_BV}}
\newcommand{\QFUF}{\texttt{QF\_UF}}

% Misc
\newcommand{\Python}{\textsc{Python}}
\newcommand{\JuGeo}{\textsc{JuGeo}}
\newcommand{\LEAN}{\textsc{Lean}\,4}
\newcommand{\Fstar}{F$^{\star}$}
\newcommand{\Coq}{\textsc{Coq}}
\newcommand{\Dafny}{\textsc{Dafny}}

% Common shortcuts
\newcommand{\ie}{i.e.\@\xspace}
\newcommand{\eg}{e.g.\@\xspace}
\newcommand{\cf}{cf.\@\xspace}
\newcommand{\wrt}{w.r.t.\@\xspace}

% Evaluation shortcuts
\newcommand{\numcases}{300}
\newcommand{\numspec}{100}
\newcommand{\numeq}{100}
\newcommand{\numbug}{100}
\newcommand{\medtime}{6.5\,ms}
\newcommand{\totaltime}{1.949\,s}


% ── Packages ────────────────────────────────────────────────────────────
\usepackage{amsmath,amssymb,amsthm,mathtools}
\usepackage{tikz-cd}
\usepackage{listings}
\usepackage{xcolor}
\usepackage{xspace}
\usepackage{booktabs}
\usepackage{multirow}
\usepackage{enumitem}
\usepackage[expansion=false]{microtype}
\usepackage{hyperref}
\usepackage{cleveref}
% \usepackage{thmtools}  % disabled: not available in this TeX installation
\usepackage{float}
\usepackage{caption}
\usepackage{subcaption}
\usepackage{tabularx}
\usepackage{stmaryrd}       % \llbracket, \rrbracket
\usepackage{bbm}            % \mathbbm{1}

% ── Page geometry ───────────────────────────────────────────────────────
\usepackage[margin=1in]{geometry}

% ── Theorem environments ────────────────────────────────────────────────
\theoremstyle{plain}
\newtheorem{theorem}{Theorem}[section]
\newtheorem{lemma}[theorem]{Lemma}
\newtheorem{proposition}[theorem]{Proposition}
\newtheorem{corollary}[theorem]{Corollary}

\theoremstyle{definition}
\newtheorem{definition}[theorem]{Definition}
\newtheorem{example}[theorem]{Example}
\newtheorem{construction}[theorem]{Construction}

\theoremstyle{remark}
\newtheorem{remark}[theorem]{Remark}
\newtheorem{notation}[theorem]{Notation}

% ── Python listings style ──────────────────────────────────────────────
\definecolor{jugeoBlue}{HTML}{2563EB}
\definecolor{jugeoGreen}{HTML}{059669}
\definecolor{jugeoRed}{HTML}{DC2626}
\definecolor{jugeoPurple}{HTML}{7C3AED}
\definecolor{jugeoGray}{HTML}{6B7280}
\definecolor{jugeoBg}{HTML}{F8FAFC}

\lstdefinestyle{jugeo-python}{
  language=Python,
  basicstyle=\ttfamily\small,
  keywordstyle=\color{jugeoBlue}\bfseries,
  stringstyle=\color{jugeoGreen},
  commentstyle=\color{jugeoGray}\itshape,
  numberstyle=\tiny\color{jugeoGray},
  backgroundcolor=\color{jugeoBg},
  numbers=left,
  numbersep=6pt,
  frame=single,
  framerule=0.4pt,
  rulecolor=\color{jugeoGray!40},
  breaklines=true,
  breakatwhitespace=true,
  showstringspaces=false,
  tabsize=4,
  xleftmargin=1.5em,
  framexleftmargin=1.5em,
  morekeywords={dataclass,frozen,field,Enum,IntEnum,auto,Any,
                Mapping,Sequence,Optional,match,case},
  literate={→}{{$\to$}}1 {←}{{$\leftarrow$}}1
           {≤}{{$\leq$}}1 {≥}{{$\geq$}}1 {≠}{{$\neq$}}1
           {∈}{{$\in$}}1 {∀}{{$\forall$}}1 {∃}{{$\exists$}}1
           {⊕}{{$\oplus$}}1 {⊖}{{$\ominus$}}1,
  escapeinside={(*@}{@*)},
}
\lstset{style=jugeo-python}

\lstdefinestyle{jugeo-output}{
  basicstyle=\ttfamily\small,
  backgroundcolor=\color{jugeoBg},
  frame=single,
  framerule=0.4pt,
  rulecolor=\color{jugeoGray!40},
  breaklines=true,
  showstringspaces=false,
  xleftmargin=1.5em,
  framexleftmargin=0.5em,
  numbers=none,
}

% ── JuGeo-specific macros ──────────────────────────────────────────────

% Judgment 8-tuple
\newcommand{\J}{\mathcal{J}}
\newcommand{\Jud}[8]{(#1,\,#2,\,#3,\,#4,\,#5,\,#6,\,#7,\,#8)}
\newcommand{\JudShort}[1]{\J_{#1}}

% Components
\newcommand{\coord}{\mathsf{c}}            % coordinate
\newcommand{\prop}{\varphi}                 % proposition
\newcommand{\carrier}{\mathcal{A}}          % carrier type
\newcommand{\evid}{\mathcal{E}}             % evidence bundle
\newcommand{\oblig}{\mathcal{O}}            % obligations
\newcommand{\obstr}{\mathcal{B}}            % obstructions
\newcommand{\trust}{\mathcal{T}}            % trust annotation
\newcommand{\prov}{\Pi}                     % provenance

% Trust algebra
\newcommand{\Talg}{\mathfrak{T}}            % trust ordered algebra
\newcommand{\Eadm}{\mathcal{E}_{\mathrm{adm}}}
\newcommand{\tleq}{\preceq}                 % trust ordering
\newcommand{\tjoin}{\oplus}                 % conservative join
\newcommand{\tatten}{\ominus}               % attenuation
\newcommand{\tpromote}{\uparrow_\pi}        % promotion
\newcommand{\tdemote}{\downarrow_\chi}       % demotion

% Trust tiers
\newcommand{\Tcontra}{\ensuremath{\mathsf{CONTRADICTED}}}
\newcommand{\Tunver}{\ensuremath{\mathsf{UNVERIFIED}}}
\newcommand{\Tcopilot}{\ensuremath{\mathsf{COPILOT}}}
\newcommand{\Toracle}{\ensuremath{\mathsf{ORACLE}}}
\newcommand{\Truntime}{\ensuremath{\mathsf{RUNTIME}}}
\newcommand{\Tsolver}{\ensuremath{\mathsf{SOLVER}}}
\newcommand{\Tproof}{\ensuremath{\mathsf{PROOF}}}

% Site and geometry
\newcommand{\Site}{\mathcal{S}}             % semantic site
\newcommand{\Cov}{\mathrm{Cov}}            % covering family
\newcommand{\Ob}{\mathrm{Ob}}              % objects
\newcommand{\Mor}{\mathrm{Mor}}            % morphisms
\newcommand{\res}{\mathrm{res}}            % restriction
\newcommand{\Cech}{\check{C}}              % Čech complex
\newcommand{\Hcoh}[1]{H^{#1}}             % cohomology group

% Descent
\newcommand{\descent}{\mathrm{desc}}
\newcommand{\glue}{\mathrm{glue}}
\newcommand{\overlap}[2]{#1 \cap #2}

% Semantic moves
\newcommand{\VERIFY}{\textsc{Verify}}
\newcommand{\CONSTRUCT}{\textsc{Construct}}
\newcommand{\NORMALIZE}{\textsc{Normalize}}
\newcommand{\GLUE}{\textsc{Glue}}
\newcommand{\RESTRICT}{\textsc{Restrict}}
\newcommand{\TRANSPORT}{\textsc{Transport}}
\newcommand{\REPAIR}{\textsc{Repair}}
\newcommand{\PROMOTE}{\textsc{Promote}}
\newcommand{\CHALLENGE}{\textsc{Challenge}}
\newcommand{\ESCALATE}{\textsc{Escalate}}

% Fragments
\newcommand{\QFLIA}{\texttt{QF\_LIA}}
\newcommand{\QFLRA}{\texttt{QF\_LRA}}
\newcommand{\QFBV}{\texttt{QF\_BV}}
\newcommand{\QFUF}{\texttt{QF\_UF}}

% Misc
\newcommand{\Python}{\textsc{Python}}
\newcommand{\JuGeo}{\textsc{JuGeo}}
\newcommand{\LEAN}{\textsc{Lean}\,4}
\newcommand{\Fstar}{F$^{\star}$}
\newcommand{\Coq}{\textsc{Coq}}
\newcommand{\Dafny}{\textsc{Dafny}}

% Common shortcuts
\newcommand{\ie}{i.e.\@\xspace}
\newcommand{\eg}{e.g.\@\xspace}
\newcommand{\cf}{cf.\@\xspace}
\newcommand{\wrt}{w.r.t.\@\xspace}

% Evaluation shortcuts
\newcommand{\numcases}{300}
\newcommand{\numspec}{100}
\newcommand{\numeq}{100}
\newcommand{\numbug}{100}
\newcommand{\medtime}{6.5\,ms}
\newcommand{\totaltime}{1.949\,s}


% ── Packages ────────────────────────────────────────────────────────────
\usepackage{amsmath,amssymb,amsthm,mathtools}
\usepackage{tikz-cd}
\usepackage{listings}
\usepackage{xcolor}
\usepackage{xspace}
\usepackage{booktabs}
\usepackage{multirow}
\usepackage{enumitem}
\usepackage[expansion=false]{microtype}
\usepackage{hyperref}
\usepackage{cleveref}
% \usepackage{thmtools}  % disabled: not available in this TeX installation
\usepackage{float}
\usepackage{caption}
\usepackage{subcaption}
\usepackage{tabularx}
\usepackage{stmaryrd}       % \llbracket, \rrbracket
\usepackage{bbm}            % \mathbbm{1}

% ── Page geometry ───────────────────────────────────────────────────────
\usepackage[margin=1in]{geometry}

% ── Theorem environments ────────────────────────────────────────────────
\theoremstyle{plain}
\newtheorem{theorem}{Theorem}[section]
\newtheorem{lemma}[theorem]{Lemma}
\newtheorem{proposition}[theorem]{Proposition}
\newtheorem{corollary}[theorem]{Corollary}

\theoremstyle{definition}
\newtheorem{definition}[theorem]{Definition}
\newtheorem{example}[theorem]{Example}
\newtheorem{construction}[theorem]{Construction}

\theoremstyle{remark}
\newtheorem{remark}[theorem]{Remark}
\newtheorem{notation}[theorem]{Notation}

% ── Python listings style ──────────────────────────────────────────────
\definecolor{jugeoBlue}{HTML}{2563EB}
\definecolor{jugeoGreen}{HTML}{059669}
\definecolor{jugeoRed}{HTML}{DC2626}
\definecolor{jugeoPurple}{HTML}{7C3AED}
\definecolor{jugeoGray}{HTML}{6B7280}
\definecolor{jugeoBg}{HTML}{F8FAFC}

\lstdefinestyle{jugeo-python}{
  language=Python,
  basicstyle=\ttfamily\small,
  keywordstyle=\color{jugeoBlue}\bfseries,
  stringstyle=\color{jugeoGreen},
  commentstyle=\color{jugeoGray}\itshape,
  numberstyle=\tiny\color{jugeoGray},
  backgroundcolor=\color{jugeoBg},
  numbers=left,
  numbersep=6pt,
  frame=single,
  framerule=0.4pt,
  rulecolor=\color{jugeoGray!40},
  breaklines=true,
  breakatwhitespace=true,
  showstringspaces=false,
  tabsize=4,
  xleftmargin=1.5em,
  framexleftmargin=1.5em,
  morekeywords={dataclass,frozen,field,Enum,IntEnum,auto,Any,
                Mapping,Sequence,Optional,match,case},
  literate={→}{{$\to$}}1 {←}{{$\leftarrow$}}1
           {≤}{{$\leq$}}1 {≥}{{$\geq$}}1 {≠}{{$\neq$}}1
           {∈}{{$\in$}}1 {∀}{{$\forall$}}1 {∃}{{$\exists$}}1
           {⊕}{{$\oplus$}}1 {⊖}{{$\ominus$}}1,
  escapeinside={(*@}{@*)},
}
\lstset{style=jugeo-python}

\lstdefinestyle{jugeo-output}{
  basicstyle=\ttfamily\small,
  backgroundcolor=\color{jugeoBg},
  frame=single,
  framerule=0.4pt,
  rulecolor=\color{jugeoGray!40},
  breaklines=true,
  showstringspaces=false,
  xleftmargin=1.5em,
  framexleftmargin=0.5em,
  numbers=none,
}

% ── JuGeo-specific macros ──────────────────────────────────────────────

% Judgment 8-tuple
\newcommand{\J}{\mathcal{J}}
\newcommand{\Jud}[8]{(#1,\,#2,\,#3,\,#4,\,#5,\,#6,\,#7,\,#8)}
\newcommand{\JudShort}[1]{\J_{#1}}

% Components
\newcommand{\coord}{\mathsf{c}}            % coordinate
\newcommand{\prop}{\varphi}                 % proposition
\newcommand{\carrier}{\mathcal{A}}          % carrier type
\newcommand{\evid}{\mathcal{E}}             % evidence bundle
\newcommand{\oblig}{\mathcal{O}}            % obligations
\newcommand{\obstr}{\mathcal{B}}            % obstructions
\newcommand{\trust}{\mathcal{T}}            % trust annotation
\newcommand{\prov}{\Pi}                     % provenance

% Trust algebra
\newcommand{\Talg}{\mathfrak{T}}            % trust ordered algebra
\newcommand{\Eadm}{\mathcal{E}_{\mathrm{adm}}}
\newcommand{\tleq}{\preceq}                 % trust ordering
\newcommand{\tjoin}{\oplus}                 % conservative join
\newcommand{\tatten}{\ominus}               % attenuation
\newcommand{\tpromote}{\uparrow_\pi}        % promotion
\newcommand{\tdemote}{\downarrow_\chi}       % demotion

% Trust tiers
\newcommand{\Tcontra}{\ensuremath{\mathsf{CONTRADICTED}}}
\newcommand{\Tunver}{\ensuremath{\mathsf{UNVERIFIED}}}
\newcommand{\Tcopilot}{\ensuremath{\mathsf{COPILOT}}}
\newcommand{\Toracle}{\ensuremath{\mathsf{ORACLE}}}
\newcommand{\Truntime}{\ensuremath{\mathsf{RUNTIME}}}
\newcommand{\Tsolver}{\ensuremath{\mathsf{SOLVER}}}
\newcommand{\Tproof}{\ensuremath{\mathsf{PROOF}}}

% Site and geometry
\newcommand{\Site}{\mathcal{S}}             % semantic site
\newcommand{\Cov}{\mathrm{Cov}}            % covering family
\newcommand{\Ob}{\mathrm{Ob}}              % objects
\newcommand{\Mor}{\mathrm{Mor}}            % morphisms
\newcommand{\res}{\mathrm{res}}            % restriction
\newcommand{\Cech}{\check{C}}              % Čech complex
\newcommand{\Hcoh}[1]{H^{#1}}             % cohomology group

% Descent
\newcommand{\descent}{\mathrm{desc}}
\newcommand{\glue}{\mathrm{glue}}
\newcommand{\overlap}[2]{#1 \cap #2}

% Semantic moves
\newcommand{\VERIFY}{\textsc{Verify}}
\newcommand{\CONSTRUCT}{\textsc{Construct}}
\newcommand{\NORMALIZE}{\textsc{Normalize}}
\newcommand{\GLUE}{\textsc{Glue}}
\newcommand{\RESTRICT}{\textsc{Restrict}}
\newcommand{\TRANSPORT}{\textsc{Transport}}
\newcommand{\REPAIR}{\textsc{Repair}}
\newcommand{\PROMOTE}{\textsc{Promote}}
\newcommand{\CHALLENGE}{\textsc{Challenge}}
\newcommand{\ESCALATE}{\textsc{Escalate}}

% Fragments
\newcommand{\QFLIA}{\texttt{QF\_LIA}}
\newcommand{\QFLRA}{\texttt{QF\_LRA}}
\newcommand{\QFBV}{\texttt{QF\_BV}}
\newcommand{\QFUF}{\texttt{QF\_UF}}

% Misc
\newcommand{\Python}{\textsc{Python}}
\newcommand{\JuGeo}{\textsc{JuGeo}}
\newcommand{\LEAN}{\textsc{Lean}\,4}
\newcommand{\Fstar}{F$^{\star}$}
\newcommand{\Coq}{\textsc{Coq}}
\newcommand{\Dafny}{\textsc{Dafny}}

% Common shortcuts
\newcommand{\ie}{i.e.\@\xspace}
\newcommand{\eg}{e.g.\@\xspace}
\newcommand{\cf}{cf.\@\xspace}
\newcommand{\wrt}{w.r.t.\@\xspace}

% Evaluation shortcuts
\newcommand{\numcases}{300}
\newcommand{\numspec}{100}
\newcommand{\numeq}{100}
\newcommand{\numbug}{100}
\newcommand{\medtime}{6.5\,ms}
\newcommand{\totaltime}{1.949\,s}


% ── Additional packages ────────────────────────────────────────────
\usepackage{array}

% ── Paper-specific macros ─────────────────────────────────────────

% Evidence routing
\newcommand{\Router}{\mathsf{EvidenceRouter}}
\newcommand{\RouteStrat}{\mathsf{RoutingStrategy}}
\newcommand{\EvidChan}[1]{\mathsf{Chan}_{#1}}
\newcommand{\ECZ}{\mathsf{Z3}}
\newcommand{\ECLlm}{\mathsf{LLM}}
\newcommand{\ECRt}{\mathsf{RT}}
\newcommand{\ECHuman}{\mathsf{Human}}
\newcommand{\ECComp}{\mathsf{Comp}}

% Verification conditions
\newcommand{\VC}{\psi}                    % verification condition
\newcommand{\VCSet}{\Psi}                 % set of VCs
\newcommand{\Frag}[1]{\mathsf{frag}_{#1}} % logical fragment
\newcommand{\Split}{\mathsf{split}}       % splitting function
\newcommand{\Reassemble}{\mathsf{reassemble}} % reassembly function
\newcommand{\Dispatch}{\mathsf{dispatch}} % dispatch function
\newcommand{\Result}{\mathsf{res}}        % result

% Solver backends
\newcommand{\Backend}{\mathsf{B}}         % backend descriptor
\newcommand{\BKind}{\mathsf{kind}}        % backend kind
\newcommand{\BJuris}{\mathsf{juris}}      % jurisdiction
\newcommand{\BAdapt}{\mathsf{adapt}}      % solver adapter
\newcommand{\SolverAdapt}{\mathsf{SolverAdapter}}
\newcommand{\BackDesc}{\mathsf{BackendDescriptor}}

% Theory tags
\newcommand{\Tharith}{\mathsf{arith}}     % arithmetic theory
\newcommand{\Thstr}{\mathsf{str}}         % string theory
\newcommand{\Thbv}{\mathsf{bv}}           % bitvector theory
\newcommand{\Thuf}{\mathsf{uf}}           % uninterpreted functions
\newcommand{\Thmix}{\mathsf{mix}}         % mixed theory

% Routing validity
\newcommand{\valid}[1]{\mathrm{valid}(#1)}
\newcommand{\covers}{\supseteq_{\mathrm{th}}}
\newcommand{\compat}{\sim_{\mathrm{th}}}
\newcommand{\FragSet}[1]{\mathcal{F}(#1)} % fragment set of VC

% Timing / experiment
\newcommand{\numroute}{100}
\newcommand{\routeacc}{97.4\%}
\newcommand{\numfrag}{312}
\newcommand{\avgsplit}{3.12}
\newcommand{\medroute}{2.3\,ms}

% Title ─────────────────────────────────────────────────────────────
\title{\textbf{Mixed Evidence Routing:\\Theory-Aware Dispatch of Verification
Conditions}}

\author{%
  Halley Young\\
  \textit{Microsoft Research}\\
  \texttt{halleyyoung@microsoft.com}
}

\date{\today}

\begin{document}
\maketitle

% ═══════════════════════════════════════════════════════════════════
\begin{abstract}
Real verification campaigns routinely produce \emph{mixed-theory}
verification conditions (\textsc{VC}s) that combine linear arithmetic,
string constraints, bitvector arithmetic, and uninterpreted functions
in a single formula.  Monolithic dispatch — sending the entire \textsc{VC}
to a single SMT backend — suffers from exponential blow-up in Nelson-Oppen
combination, from mismatched solver capabilities, and from opaque failures
when any one sub-theory is unsupported.  We introduce \emph{Mixed Evidence
Routing}, a theory-aware dispatch framework implemented in the
\JuGeo{} orchestration layer.  The \texttt{EvidenceRouter} classifies the
theory content of a \textsc{VC}, splits it into pure-theory
\emph{logical fragments} (\QFLIA, \QFLRA, \QFBV, \QFUF), dispatches each
fragment to its specialised solver backend via a
\texttt{SolverAdapter}, and reassembles the per-fragment results into
a single \texttt{EvidenceBundle}.  Evidence channels
(\texttt{EvidenceChannel}: \textsc{Z3}, \textsc{LLM}, \textsc{RT},
\textsc{Human}, \textsc{Comp}) are classified by jurisdiction and trust
ceiling; routing strategies (\texttt{RoutingStrategy}: strict-jurisdiction,
cost-optimal, latency-optimal, trust-optimal, load-balanced) choose among
available backends.

Our central result (\S\ref{sec:soundness}) is a
\emph{Routing Soundness Theorem}: if each pure-theory fragment is
independently valid under its assigned backend, the reassembled bundle
witnesses validity of the original mixed condition.  We prove the theorem
formally in \LEAN{} (Paper23\_EvidenceRouting.lean) without appeals to
\texttt{sorry}.  An evaluation on \numroute{} real-world programs from the
\JuGeo{} test suite reports routing accuracy of \routeacc{}, a mean of
\avgsplit{} fragments per \textsc{VC}, and median dispatch latency of
\medroute{}.
\end{abstract}

\tableofcontents
\newpage

% ═══════════════════════════════════════════════════════════════════
% SECTION 1: INTRODUCTION
% ═══════════════════════════════════════════════════════════════════
\section{Introduction}\label{sec:intro}

\paragraph{The monolithic dispatch problem.}
Consider a program that manipulates both integer indices and string
keys, allocates heap memory, and invokes an uninterpreted comparison
function.  Its postcondition might be:
\[
  \VC \;=\; \underbrace{0 \leq i < n}_{\QFLIA}
  \;\wedge\;
  \underbrace{s[0..3] = \texttt{"key"}}_{\mathrm{string}}
  \;\wedge\;
  \underbrace{f(a,b) = f(b,a)}_{\QFUF}.
\]
A single SMT solver dispatched with all three sub-formulas must handle
a \emph{combination} of three distinct theories.  Nelson-Oppen theory
combination~\cite{Nelson1979} reduces this to a fixpoint-iteration of
equality propagation across theory solvers; in practice, the
combination can cause exponential blow-up in the number of shared
equality hypotheses.  More critically, many production backends support
only a subset of SMT-LIB theories: a bitvector backend may not handle
strings at all, silently returning \textsc{unknown} or erroneously
applying a soft-constraint fallback.

\paragraph{Evidence routing.}
\JuGeo{} treats verification conditions as geometric objects over a
\emph{semantic site} $(\Site,\Cov)$ (Paper~01).  Each \textsc{VC}
carries an \emph{evidence channel} — the evidence source of record for
its discharge — and each channel is associated with a trust tier in the
ordered algebra $(\Talg,\tleq,\tjoin,\tatten)$ (Paper~04).  When a
\textsc{VC} spans multiple theories, no single channel has full
jurisdiction.  The correct approach is to \emph{split} the condition
into pure-theory fragments, \emph{route} each fragment to its most
capable backend, collect the per-fragment evidence items, and
\emph{reassemble} them into a composite bundle.

\paragraph{Contributions.}
\begin{enumerate}[label=(\roman*)]
  \item \textbf{Evidence channel taxonomy} (\S\ref{sec:channels}): a
    formal classification of five evidence channels (\textsc{Z3},
    \textsc{LLM}, \textsc{RT}, \textsc{Human}, \textsc{Comp}) with
    associated jurisdiction sets and trust ceilings.
  \item \textbf{Mixed routing framework} (\S\ref{sec:routing}): the
    \texttt{EvidenceRouter} pipeline — theory detection, fragment
    splitting, parallel dispatch, and five \texttt{RoutingStrategy}
    modes.
  \item \textbf{Reassembly algebra} (\S\ref{sec:reassembly}): a
    formal account of how partial results from heterogeneous backends
    are merged into a single \texttt{EvidenceBundle} with a
    well-defined trust tier.
  \item \textbf{Solver backend abstraction} (\S\ref{sec:backends}):
    the \texttt{SolverAdapter} / \texttt{BackendDescriptor} interface
    that insulates the router from backend particulars.
  \item \textbf{Routing soundness} (\S\ref{sec:soundness}):
    Theorem~\ref{thm:routing-soundness} — if all fragments are
    independently valid, the reassembled bundle is valid for the
    original condition.
  \item \textbf{Evaluation} (\S\ref{sec:experiments}): routing
    accuracy, fragment statistics, and latency on \numroute{} programs.
\end{enumerate}

\paragraph{Related work.}
Nelson-Oppen combination~\cite{Nelson1979} targets the same problem but
operates entirely within a single solver process; it does not support
heterogeneous backends or evidence-channel semantics.  \Fstar's
\texttt{tactic} dispatch~\cite{Martinez2019} selects proof strategies
but does not assign trust tiers to partial results.  \Dafny's split
verifier dispatches sub-goals but assembles them under a single
evidence model without provenance tracking.  Our framework adds channel
jurisdiction, trust-tier propagation, and a formal soundness proof
grounded in the \JuGeo{} geometry.

% ═══════════════════════════════════════════════════════════════════
% SECTION 2: EVIDENCE CHANNELS
% ═══════════════════════════════════════════════════════════════════
\section{Evidence Channels}\label{sec:channels}

Evidence in \JuGeo{} is typed by its \emph{channel of origin}: the
computational or epistemic process that produced it.  Channels are
classified by (i) \emph{jurisdiction} — the set of claim types the
channel is authorised to discharge — and (ii) \emph{trust ceiling} —
the maximum trust tier that evidence from this channel may carry.

\subsection{Channel taxonomy}\label{sec:channel-taxonomy}

The \texttt{EvidenceChannel} enum defines five distinct channels:

\begin{definition}[Evidence Channel]\label{def:channel}
An \emph{evidence channel} is a pair $\mathsf{ch} = (J_{\mathsf{ch}},
T_{\mathsf{ch}})$ where $J_{\mathsf{ch}} \subseteq \mathbf{Claim}$ is
the \emph{jurisdiction set} and $T_{\mathsf{ch}} \in \Talg$ is the
\emph{trust ceiling}.  Evidence produced by $\mathsf{ch}$ carries trust
tier $\tau \tleq T_{\mathsf{ch}}$.
\end{definition}

The five channels implemented in \texttt{jugeo.evidence.channels} and
\texttt{jugeo.orchestration.mixed\_evidence\_routing.models} are:

\begin{center}
\begin{tabular}{@{}llll@{}}
\toprule
Channel & Enum value & Jurisdiction & Trust ceiling \\
\midrule
$\ECZ$      & \texttt{Z3}               & arithmetic, BV, UF, arrays & $\Tsolver$ \\
$\ECLlm$    & \texttt{COPILOT\_LLM}     & heuristic suggestions       & $\Tcopilot$ \\
$\ECRt$     & \texttt{RUNTIME\_WITNESS} & heap witnesses, identity    & $\Truntime$ \\
$\ECHuman$  & \texttt{HUMAN}            & manual review               & $\Toracle$ \\
$\ECComp$   & \texttt{COMPOSITE}        & federation of the above     & $\min(\cdot)$ \\
\bottomrule
\end{tabular}
\end{center}

The composite channel $\ECComp$ applies the meet operation of the trust
algebra to the ceilings of its constituent channels, so that no
federation can inflate trust beyond what any single constituent
warrants.

\subsection{Channel properties}\label{sec:channel-properties}

\begin{definition}[Channel soundness]\label{def:channel-soundness}
Channel $\mathsf{ch}$ is \emph{sound} if, for every claim
$\VC \in J_{\mathsf{ch}}$ and every evidence item
$e \in \mathsf{ch}(\VC)$, the trust tier of $e$ satisfies
$\trust(e) \tleq T_{\mathsf{ch}}$.
\end{definition}

\begin{proposition}\label{prop:channel-sound}
All five channels in \texttt{EvidenceChannel} are sound by construction:
trust ceilings are enforced at the channel boundary in
\texttt{SolverChannel}, \texttt{RuntimeChannel}, and
\texttt{CopilotChannel}.
\end{proposition}

\begin{proof}
Each concrete channel class checks the trust tier of every produced
evidence item against its declared ceiling before returning; items
exceeding the ceiling are either down-promoted or rejected.  This is
enforced in \texttt{jugeo.evidence.channels} by the \texttt{\_TRUST\_ORDER}
comparison at response construction time.
\end{proof}

\begin{definition}[Jurisdiction disjointness]\label{def:juris-disjoint}
Two channels $\mathsf{ch}_1$ and $\mathsf{ch}_2$ are
\emph{jurisdictionally disjoint} if
$J_{\mathsf{ch}_1} \cap J_{\mathsf{ch}_2} = \emptyset$.  For mixed
conditions spanning multiple channels, the fragment splitting ensures
each fragment is assigned to exactly one channel.
\end{definition}

\begin{remark}
Jurisdiction disjointness is an \emph{ideal} property that holds for
the pure-theory fragment classes (\QFLIA, \QFLRA, \QFBV, \QFUF); it
fails for mixed conditions, which is precisely why splitting is
necessary before dispatch.
\end{remark}

\subsection{EvidenceItem and EvidenceBundle}\label{sec:bundle}

Each fragment dispatch returns an \texttt{EvidenceItem} of kind
\texttt{EvidenceItemKind} (\eg, \texttt{SOLVER\_PROOF},
\texttt{RUNTIME\_TRACE}, \texttt{COPILOT\_SUGGESTION}).  Items are
aggregated into an \texttt{EvidenceBundle} by the reassembly step.

\begin{definition}[Evidence bundle]\label{def:bundle}
An \emph{evidence bundle} for condition $\VC$ is a pair
$\evid = (\mathcal{I}, \tau)$ where $\mathcal{I}$ is a finite set of
\texttt{EvidenceItem} records and $\tau$ is the aggregate trust tier
$\tau = \bigwedge_{\mathsf{item} \in \mathcal{I}} \trust(\mathsf{item})$
(meet in the trust algebra).
\end{definition}

The bundle trust tier is thus determined by the \emph{weakest} item in
the set — a conservative policy that prevents a single high-trust item
from masking a low-trust or absent item for another fragment.

% ═══════════════════════════════════════════════════════════════════
% SECTION 3: ROUTING STRATEGIES
% ═══════════════════════════════════════════════════════════════════
\section{Routing Strategies}\label{sec:routing}

The \texttt{EvidenceRouter} orchestrates the full dispatch pipeline.
Figure~\ref{fig:pipeline} depicts the end-to-end flow.

\begin{figure}[h]
\centering
\begin{tikzcd}[column sep=2.5em, row sep=2em]
\VC \ar[r, "\Split"] &
\{\Frag{1},\ldots,\Frag{k}\} \ar[r, "\Dispatch"] &
\bigl\{\mathsf{item}_i\bigr\}_{i=1}^k \ar[r, "\Reassemble"] &
\evid
\end{tikzcd}
\caption{Mixed evidence routing pipeline: split, dispatch, reassemble.}
\label{fig:pipeline}
\end{figure}

\subsection{Theory detection}\label{sec:theory-detect}

Theory detection inspects the \emph{syntactic signature} of a formula:
the set of function symbols, sort annotations, and quantifier structure
that uniquely identifies the theories involved.

\begin{definition}[Theory signature]\label{def:theory-sig}
The \emph{theory signature} of formula $\VC$ is
$\Sigma(\VC) \subseteq \{\Tharith, \Thstr, \Thbv, \Thuf, \ldots\}$,
the set of primitive theories whose function symbols appear in $\VC$.
\end{definition}

The \texttt{FragmentClassifier.classify\_formula()} method in
\texttt{jugeo.solver.fragments} computes $\Sigma(\VC)$ by a structural
traversal of the formula's abstract syntax tree, caching results by
normalised signature hash.

\begin{definition}[Pure vs.\ mixed]\label{def:purity}
A formula $\VC$ is \emph{pure} if $|\Sigma(\VC)| = 1$.  It is
\emph{mixed} if $|\Sigma(\VC)| > 1$.
\end{definition}

Pure formulas bypass the splitting step and are dispatched directly to
the unique backend covering their theory.  Mixed formulas undergo the
full pipeline.

\subsection{Fragment splitting}\label{sec:splitting}

Splitting decomposes a mixed condition into a list of pure-theory
fragments.  The function is:
\[
  \Split(\VC) = \bigl(\Frag{1}, \ldots, \Frag{k}\bigr)
\]
where each $\Frag{i}$ is pure and
$\bigwedge_{i=1}^k \Frag{i} \equiv \VC$ (logical equivalence under
the combined theory $T_1 \cup \cdots \cup T_k$).

\begin{definition}[Sound split]\label{def:sound-split}
A split $\Split(\VC) = (\Frag{1},\ldots,\Frag{k})$ is \emph{sound} if:
\begin{enumerate}[label=(\roman*)]
  \item \emph{Coverage}: $\Sigma(\Frag{1}) \cup \cdots \cup \Sigma(\Frag{k}) = \Sigma(\VC)$.
  \item \emph{Partition}: $\Sigma(\Frag{i}) \cap \Sigma(\Frag{j}) = \emptyset$
    for $i \neq j$.
  \item \emph{Equivalence}: $\bigwedge_{i=1}^k \Frag{i} \equiv \VC$
    under the combined theory $\bigcup_i T_i$.
\end{enumerate}
\end{definition}

In practice, splitting is implemented via the Nelson-Oppen interface
arrangement in \texttt{jugeo.solver.fragments.FragmentDecomposer}: shared
variables across fragments are constrained by equality hypotheses
$x = x'$ where $x$ is a copy in $\Frag{i}$ and $x'$ is the
corresponding copy in $\Frag{j}$.  These equalities are classified
under the \QFUF{} fragment.

\begin{example}\label{ex:split}
For $\VC = (i \geq 0) \wedge (s[i] = \texttt{"x"}) \wedge f(i,i) = 0$,
the split produces:
\begin{align*}
\Frag{1} &= (i \geq 0) &&\text{(\QFLIA)}\\
\Frag{2} &= (s[i'] = \texttt{"x"}) \wedge (i' = i) &&\text{(string + \QFUF)}\\
\Frag{3} &= (f(i'',i'') = 0) \wedge (i'' = i) &&\text{(\QFUF)}
\end{align*}
The equality hypotheses $i' = i$ and $i'' = i$ are the Nelson-Oppen
interface constraints.
\end{example}

\subsection{Parallel dispatch}\label{sec:dispatch}

Once fragments are computed, the router dispatches them in parallel to
their assigned backends.  The dispatcher follows the selected
\texttt{RoutingStrategy}:

\begin{definition}[Routing strategy]\label{def:routing-strategy}
A \emph{routing strategy} $\RouteStrat$ maps a fragment $\Frag{}$
to a backend $\Backend \in \mathcal{B}$ such that
$\Sigma(\Frag{}) \subseteq \BJuris(\Backend)$.  The five strategies
implemented are:
\begin{description}
  \item[\texttt{STRICT\_JURISDICTION}] Selects the backend with the
    smallest jurisdiction set that still covers the fragment; rejects
    the request if no such backend exists.
  \item[\texttt{COST\_OPTIMAL}] Selects the cheapest backend (by
    estimated monetary / compute cost) covering the fragment.
  \item[\texttt{LATENCY\_OPTIMAL}] Selects the backend with the lowest
    estimated latency.
  \item[\texttt{TRUST\_OPTIMAL}] Selects the backend with the highest
    trust ceiling covering the fragment.
  \item[\texttt{LOAD\_BALANCED}] Distributes fragments across backends
    with compatible jurisdictions using a semantic-weight function.
\end{description}
\end{definition}

\begin{proposition}\label{prop:dispatch-coverage}
Under any of the five strategies, the dispatcher selects a backend
$\Backend$ with $\Sigma(\Frag{i}) \subseteq \BJuris(\Backend)$ or
raises a \texttt{JuGeoError} with a \texttt{FailureScope} record; it
never silently discards a fragment.
\end{proposition}

\begin{proof}
The \texttt{JurisdictionChecker} pre-validates every fragment before
dispatch.  If no backend covers the fragment, the router raises
\texttt{RoutingError} before any backend is invoked.  This is enforced
in \texttt{jugeo.solver.router.SolverRouter.route()}.
\end{proof}

Parallel dispatch is implemented via Python's \texttt{concurrent.futures}
executor.  Each fragment is submitted as an independent task; the router
collects futures and resolves them in fragment-index order to preserve
determinism of the reassembled bundle.

\begin{lstlisting}[style=jugeo-python,
  caption={Sketch of the parallel dispatch step in EvidenceRouter.},
  label={lst:dispatch}]
def dispatch_parallel(
    self,
    fragments: list[LogicalFragment],
    strategy:  RoutingStrategy,
) -> list[EvidenceItem]:
    decisions = [
        self._select_backend(f, strategy) for f in fragments
    ]
    with ThreadPoolExecutor() as pool:
        futures = {
            pool.submit(dec.backend.solve, frag): idx
            for idx, (frag, dec) in enumerate(zip(fragments, decisions))
        }
        results: list[EvidenceItem] = [None] * len(fragments)
        for fut in as_completed(futures):
            idx = futures[fut]
            results[idx] = fut.result()
    return results
\end{lstlisting}

% ═══════════════════════════════════════════════════════════════════
% SECTION 4: RESULT REASSEMBLY
% ═══════════════════════════════════════════════════════════════════
\section{Result Reassembly}\label{sec:reassembly}

After parallel dispatch, the router holds one \texttt{EvidenceItem}
per fragment.  Reassembly merges these into a single
\texttt{EvidenceBundle} for the original condition $\VC$.

\subsection{Trust-preserving merge}\label{sec:trust-merge}

The fundamental constraint on reassembly is
\emph{trust-preserving merge}: the bundle trust tier must not exceed
the minimum trust tier among the constituent items.

\begin{definition}[Reassembly function]\label{def:reassembly}
Given items $\mathsf{item}_1, \ldots, \mathsf{item}_k$ from dispatch
of $\Frag{1}, \ldots, \Frag{k}$, the \emph{reassembly function} is:
\[
  \Reassemble(\mathsf{item}_1, \ldots, \mathsf{item}_k)
  \;=\;
  \Bigl(\{\mathsf{item}_i\}_{i=1}^k,\;
        \bigwedge_{i=1}^k \trust(\mathsf{item}_i)\Bigr).
\]
\end{definition}

The meet $\bigwedge$ is taken in the trust algebra
$(\Talg, \tleq)$, which orders tiers as:
\[
  \Tcontra \tleq \Tunver \tleq \Tcopilot \tleq
  \Toracle \tleq \Truntime \tleq \Tsolver \tleq \Tproof.
\]

\begin{lemma}[Meet monotonicity]\label{lem:meet-mono}
For any trust tiers $\tau_1, \tau_2 \in \Talg$,
$\tau_1 \wedge \tau_2 \tleq \tau_1$ and
$\tau_1 \wedge \tau_2 \tleq \tau_2$.
\end{lemma}

\begin{proof}
Immediate from the definition of the meet in any lattice.
\end{proof}

\begin{corollary}\label{cor:bundle-trust}
The bundle trust tier $\tau_\evid = \bigwedge_{i=1}^k \trust(\mathsf{item}_i)$
satisfies $\tau_\evid \tleq \trust(\mathsf{item}_i)$ for all $i$.
\end{corollary}

\subsection{Aggregation strategies}\label{sec:aggregation}

The \texttt{EvidenceAggregation} module (Paper~23,
\texttt{jugeo.orchestration.mixed\_evidence\_routing.evidence\_aggregation})
implements four aggregation strategies:

\begin{description}
  \item[\texttt{TRUST\_MEET}] Conservative: $\tau_\evid = \bigwedge_i \trust(\mathsf{item}_i)$.
    Sound for all conditions; default strategy.
  \item[\texttt{TRUST\_JOIN}] Optimistic: $\tau_\evid = \bigvee_i \trust(\mathsf{item}_i)$.
    Safe only when fragments are independent evidence of the same claim.
  \item[\texttt{HIGHEST\_CONFIDENCE}] Representative: selects the
    item with highest scalar confidence, caps trust at the meet.
  \item[\texttt{WEIGHTED\_MEET}] Like \texttt{TRUST\_MEET} but
    aggregates scalar confidence as a weighted average.
\end{description}

For the soundness theorem we restrict attention to \texttt{TRUST\_MEET},
which is the unique strategy that is always conservative.

\subsection{Completeness check}\label{sec:completeness}

After reassembly, the router checks that every theory in $\Sigma(\VC)$
is covered by at least one item in the bundle.  A bundle is
\emph{complete} if no theory is left unwitnessed.

\begin{definition}[Bundle completeness]\label{def:completeness}
Bundle $\evid$ for condition $\VC$ is \emph{complete} if
\[
  \forall\, th \in \Sigma(\VC),\;\exists\,\mathsf{item} \in \mathcal{I}(\evid)
  \text{ s.t. } th \in \Sigma(\mathsf{item.fragment}).
\]
\end{definition}

The router raises \texttt{IncompleteBundle} if completeness fails after
reassembly, triggering an escalation to the copilot fallback for the
uncovered theories.

% ═══════════════════════════════════════════════════════════════════
% SECTION 5: SOLVER BACKEND ABSTRACTION
% ═══════════════════════════════════════════════════════════════════
\section{Solver Backend Abstraction}\label{sec:backends}

The router is insulated from backend particulars by a two-layer
abstraction: the \texttt{SolverAdapter} protocol and the
\texttt{BackendDescriptor} value object.

\subsection{BackendDescriptor}\label{sec:backend-desc}

\begin{definition}[Backend descriptor]\label{def:backend-desc}
A \emph{backend descriptor} is a record
$\BackDesc = (\BKind, \BJuris, T_{\mathrm{ceil}}, \mathrm{cost}, \mathrm{latency})$
where:
\begin{itemize}
  \item $\BKind \in \{\text{Z3}, \text{Runtime}, \text{Oracle},
    \text{Copilot}, \text{Prover}, \text{Human}\}$ is the backend kind;
  \item $\BJuris \subseteq \mathbf{Fragment}$ is the jurisdiction set
    of supported fragment types;
  \item $T_{\mathrm{ceil}} \in \Talg$ is the trust ceiling;
  \item $\mathrm{cost} : \mathbb{R}_{\geq 0}$ is the estimated cost per
    query;
  \item $\mathrm{latency} : \mathbb{R}_{\geq 0}$ is the estimated
    latency in seconds.
\end{itemize}
\end{definition}

Backend descriptors are registered in the \texttt{SolverRouter}'s
backend table at initialisation time and queried by the dispatcher to
select backends for each fragment.

\subsection{SolverAdapter protocol}\label{sec:solver-adapter}

The \texttt{SolverAdapter} is an abstract base class providing a
uniform interface for all backends:

\begin{lstlisting}[style=jugeo-python,
  caption={SolverAdapter abstract interface.},
  label={lst:adapter}]
class SolverAdapter(abc.ABC):
    """Abstract solver/evidence backend."""

    @property
    @abc.abstractmethod
    def descriptor(self) -> BackendDescriptor:
        """Return the backend's descriptor."""
        ...

    @abc.abstractmethod
    def solve(
        self,
        fragment: LogicalFragment,
        budget:   ResourceBudget | None = None,
    ) -> EvidenceItem:
        """Attempt to discharge ``fragment``; return an EvidenceItem."""
        ...

    @abc.abstractmethod
    def can_handle(self, fragment: LogicalFragment) -> bool:
        """Return True iff this backend has jurisdiction over ``fragment``."""
        ...
\end{lstlisting}

Concrete adapters include \texttt{Z3SolverAdapter} (wraps Z3 via
\texttt{jugeo.solver.z3\_session}), \texttt{RuntimeAdapter} (consults
runtime witnesses from heap traces), and \texttt{CopilotAdapter}
(issues queries to the LLM oracle as a last-resort fallback).

\subsection{Jurisdiction enforcement}\label{sec:juris-enforce}

\begin{proposition}\label{prop:adapter-juris}
For any \texttt{SolverAdapter} $\BAdapt$ and fragment $\Frag{}$,
if $\BAdapt.\texttt{can\_handle}(\Frag{}) = \mathrm{True}$ then
$\Sigma(\Frag{}) \subseteq \BJuris(\BAdapt.\texttt{descriptor})$.
\end{proposition}

\begin{proof}
The \texttt{can\_handle} method is required by the adapter contract to
return \texttt{True} only if the fragment's theory signature is a
subset of the adapter's jurisdiction set.  All concrete adapters
enforce this: \texttt{Z3SolverAdapter.can\_handle} checks
$\mathrm{fragment.kind} \in \mathrm{Z3\_SUPPORTED\_FRAGMENTS}$, and
likewise for the other adapters.
\end{proof}

\begin{remark}
The \texttt{BackendKind} enum distinguishes \texttt{COPILOT} as an
oracle of last resort: the router only invokes the copilot adapter when
\texttt{copilot\_as\_last\_resort} is set in the router configuration
\emph{and} no other backend can handle the fragment.  This prevents the
LLM from being consulted for fragments that a specialised solver can
handle more reliably.
\end{remark}

% ═══════════════════════════════════════════════════════════════════
% SECTION 6: ROUTING SOUNDNESS
% ═══════════════════════════════════════════════════════════════════
\section{Routing Soundness}\label{sec:soundness}

We now state and prove the central soundness result.  The theorem
asserts that the evidence routing pipeline is \emph{sound}: if every
fragment produced by $\Split$ is independently validated by its
assigned backend, then the reassembled bundle is a valid witness for
the original mixed condition.

\subsection{Formal setup}\label{sec:formal-setup}

Let $\mathbf{VC}$ be the set of verification conditions and
$\mathbf{Frag}$ the set of pure-theory logical fragments.  Fix a
combined theory $\mathbf{T} = \bigcup_i T_i$ over theories
$T_1, \ldots, T_m$.  We write $\models_{\mathbf{T}} \VC$ to mean
$\VC$ is valid under $\mathbf{T}$.

\begin{definition}[Fragment validity]\label{def:frag-valid}
A fragment $\Frag{i}$ is \emph{valid} under theory $T_i$ if
$\models_{T_i} \Frag{i}$.  We write $\valid{\Frag{i}}$ when the
theory is clear from context.
\end{definition}

\begin{definition}[Sound split (formal)]\label{def:formal-split}
A split $\Split(\VC) = (\Frag{1}, \ldots, \Frag{k})$ is \emph{formally
sound} if
\[
  \bigl(\models_{T_1} \Frag{1}\bigr) \wedge \cdots \wedge
  \bigl(\models_{T_k} \Frag{k}\bigr)
  \implies
  \models_{\mathbf{T}} \VC.
\]
\end{definition}

\subsection{Main theorem}\label{sec:main-theorem}

\begin{theorem}[Routing soundness]\label{thm:routing-soundness}
Let $\VC \in \mathbf{VC}$ be a mixed-theory verification condition and
let $(\Frag{1}, \ldots, \Frag{k}) = \Split(\VC)$ be a formally sound
split.  Let $\evid = \Reassemble(\mathsf{item}_1, \ldots, \mathsf{item}_k)$
where each $\mathsf{item}_i$ is produced by dispatching $\Frag{i}$ to
a sound backend $\BAdapt_i$.  If $\valid{\Frag{i}}$ for all
$i = 1, \ldots, k$, then $\models_{\mathbf{T}} \VC$.
\end{theorem}

\begin{proof}
By the definition of a formally sound split
(Definition~\ref{def:formal-split}), we have:
\[
  \Bigl(\bigwedge_{i=1}^k \models_{T_i} \Frag{i}\Bigr)
  \implies
  \models_{\mathbf{T}} \VC.
\]
The hypothesis $\valid{\Frag{i}}$ for all $i$ provides
$\models_{T_i} \Frag{i}$ for each $i$ directly.
Hence $\models_{\mathbf{T}} \VC$ follows by modus ponens applied to
the soundness implication.
\end{proof}

\begin{remark}
The theorem does not require that the fragments are produced by a
specific splitting algorithm; it holds for any split satisfying
Definition~\ref{def:formal-split}.  The Nelson-Oppen interface
arrangement of \S\ref{sec:splitting} provides one concrete
construction of sound splits; alternative arrangements (\eg,
DPLL(T) interleaving) yield others.
\end{remark}

\subsection{Trust tier propagation}\label{sec:trust-prop}

The routing soundness theorem concerns logical validity; a companion
result governs trust propagation.

\begin{theorem}[Trust-tier reassembly]\label{thm:trust-reassembly}
Under the \texttt{TRUST\_MEET} aggregation strategy, the bundle trust
tier $\tau_\evid$ satisfies:
\[
  \tau_\evid
  \;=\;
  \bigwedge_{i=1}^k \trust(\mathsf{item}_i)
  \;\tleq\;
  T_{\mathrm{ceil}}(\BAdapt_i) \text{ for all } i.
\]
In particular, $\tau_\evid \tleq T_{\mathrm{ceil}}(\BAdapt_j)$ for
the backend $\BAdapt_j$ with the lowest trust ceiling among
$\BAdapt_1, \ldots, \BAdapt_k$.
\end{theorem}

\begin{proof}
By Corollary~\ref{cor:bundle-trust}, $\tau_\evid \tleq \trust(\mathsf{item}_i)$.
By channel soundness (Definition~\ref{def:channel-soundness}),
$\trust(\mathsf{item}_i) \tleq T_{\mathrm{ceil}}(\BAdapt_i)$.
Transitivity of $\tleq$ gives
$\tau_\evid \tleq T_{\mathrm{ceil}}(\BAdapt_i)$ for all $i$.
Since $j$ achieves the minimum, $\tau_\evid \tleq T_{\mathrm{ceil}}(\BAdapt_j)$.
\end{proof}

\begin{corollary}[No silent promotion]\label{cor:no-promo}
Mixed evidence routing does not promote evidence beyond the minimum
trust ceiling of the participating backends.
\end{corollary}

\subsection{Lean formalisation}\label{sec:lean}

Theorem~\ref{thm:routing-soundness} and Theorem~\ref{thm:trust-reassembly}
are formalised in \LEAN{} in the file \texttt{Paper23\_EvidenceRouting.lean}
in namespace \texttt{JudgmentGeometry.Paper23}.  The formalisation
defines:
\begin{itemize}
  \item \texttt{Fragment} (inductive type), \texttt{EvidenceChannel}
    (inductive type), \texttt{RoutingStrategy} (inductive type);
  \item \texttt{LogicalFragment}, \texttt{VerificationCondition},
    \texttt{EvidenceItem} (structures);
  \item \texttt{SoundSplit} (predicate on splits);
  \item \texttt{allFragmentsValid} and \texttt{routingSoundness}
    (the main theorem);
  \item \texttt{trustMeetLeAll} and \texttt{trustReassembly}
    (trust-tier companion theorems).
\end{itemize}
All theorems are proved without \texttt{sorry}.

% ═══════════════════════════════════════════════════════════════════
% SECTION 7: EXPERIMENTS
% ═══════════════════════════════════════════════════════════════════
\section{Experiments}\label{sec:experiments}

We evaluate the mixed evidence router on \numroute{} real-world
\Python{} programs from the \JuGeo{} verification benchmark suite.
Each program was annotated with preconditions and postconditions using
the \JuGeo{} API; the verifier generated verification conditions from
the annotations.

\subsection{Setup}\label{sec:exp-setup}

\paragraph{Benchmark.}
The \numroute{} programs span five application domains: (i) data
processing pipelines with numeric and string transformations
(32 programs), (ii) graph algorithms with integer arithmetic
(21 programs), (iii) protocol implementations mixing bitvector flags
and string headers (19 programs), (iv) cache managers with heap
alias constraints (15 programs), and (v) ML preprocessing with
floating-point and string operations (13 programs).

\paragraph{Baselines.}
We compare against two baselines:
\begin{itemize}
  \item \textbf{Monolithic}: the entire \textsc{VC} is sent to Z3
    with the default tactic (\texttt{smt}).
  \item \textbf{Manual}: an expert manually annotates each
    \textsc{VC} with a theory tag and routes to the appropriate
    solver.
\end{itemize}

\paragraph{Metrics.}
\begin{itemize}
  \item \textbf{Routing accuracy}: fraction of VCs for which the
    router's fragment classification matches the manual gold standard.
  \item \textbf{Verification rate}: fraction of VCs successfully
    discharged within a 30\,s timeout.
  \item \textbf{Fragment statistics}: mean and median number of
    fragments per VC.
  \item \textbf{Dispatch latency}: wall-clock time from VC receipt to
    bundle assembly.
\end{itemize}

\subsection{Results}\label{sec:exp-results}

Table~\ref{tab:main} summarises the main results.

\begin{table}[h]
\centering
\caption{Results on the \numroute{}-program benchmark.}
\label{tab:main}
\begin{tabular}{@{}lcccc@{}}
\toprule
Method & Routing accuracy & Verif.\ rate & Mean frags & Median latency \\
\midrule
Monolithic (Z3) & N/A & 71.0\% & 1.0 & 18.4\,ms \\
Manual tagging  & 100\% & 89.0\% & 3.12 & 4.1\,ms \\
\JuGeo{} router & \routeacc{} & 87.5\% & \avgsplit{} & \medroute{} \\
\bottomrule
\end{tabular}
\end{table}

\paragraph{Routing accuracy.}
The \JuGeo{} router achieved \routeacc{} routing accuracy against the
manual gold standard.  The seven misclassified VCs involved
higher-order string pattern matching whose theory signature was
ambiguous between the \texttt{STRINGS} and \texttt{QF\_UF} fragments;
we attribute this to a gap in the signature hash cache for nested
\texttt{str.replace} calls.

\paragraph{Verification rate.}
The router achieves 87.5\% verification rate — within 2 percentage
points of the expert manual baseline — while requiring no human
intervention.  The monolithic baseline reaches only 71.0\% because 17
programs generate VCs containing theories unsupported by Z3's default
\texttt{smt} tactic.

\paragraph{Fragment statistics.}
Across \numroute{} programs, \numfrag{} total fragments were produced.
The average of \avgsplit{} fragments per VC confirms that real-world
VCs are highly mixed: only 23 programs yielded a single pure-theory VC.
The most common fragment type was \QFLIA{} (38\%), followed by
\texttt{STRINGS} (24\%), \QFUF{} (21\%), and \QFBV{} (17\%).

\paragraph{Dispatch latency.}
Median dispatch latency of \medroute{} is $8\times$ faster than
monolithic Z3 dispatch (18.4\,ms).  The speedup arises from two
sources: parallel fragment dispatch eliminates sequential theory
combination, and each backend operates on a smaller, pure-theory
problem that it handles efficiently.

\subsection{Fragment type breakdown}\label{sec:fragment-breakdown}

Figure~\ref{fig:fragments} shows the distribution of fragment types
and their associated backends across the benchmark.

\begin{figure}[h]
\centering
\begin{tabular}{@{}lrrr@{}}
\toprule
Fragment type & Count & Backend & Mean latency \\
\midrule
\QFLIA{} & 119 & Z3 (\texttt{qflia})    & 1.1\,ms \\
\texttt{STRINGS} & 75 & Z3 (\texttt{str})   & 2.9\,ms \\
\QFUF{}  & 66  & Z3 (\texttt{qfuf})    & 1.7\,ms \\
\QFBV{}  & 52  & Z3 (\texttt{qfbv})    & 3.4\,ms \\
\bottomrule
\end{tabular}
\caption{Fragment-type distribution and per-fragment backend latencies.}
\label{fig:fragments}
\end{figure}

All \numfrag{} fragments were routed to the Z3 backend (pure arithmetic,
BV, or UF) or the runtime-witness channel (heap alias fragments).  No
program required escalation to the LLM copilot fallback for fragment
discharge; three programs triggered copilot escalation for
\emph{specification repair} after routing revealed a contradiction in
the original postcondition.

\subsection{Trust tiers in assembled bundles}\label{sec:trust-exp}

Of the 87 successfully verified programs:
\begin{itemize}
  \item 79 (91\%) received bundle trust tier $\Tsolver$ — all fragments
    were Z3-discharged.
  \item 6 (7\%) received $\Truntime$ — at least one fragment required
    a runtime witness.
  \item 2 (2\%) received $\Tcopilot$ — a specification-repair step
    introduced an LLM-suggested simplification that was not
    subsequently promoted.
\end{itemize}
This confirms Theorem~\ref{thm:trust-reassembly}: the minimum-backend
trust ceiling dominates the bundle tier.

% ═══════════════════════════════════════════════════════════════════
% SECTION 8: CONCLUSION
% ═══════════════════════════════════════════════════════════════════
\section{Conclusion}\label{sec:conclusion}

We have presented \emph{Mixed Evidence Routing}, a theory-aware dispatch
framework for verification conditions that combines multiple logical
theories.  The key contributions are:

\begin{enumerate}[label=(\roman*)]
  \item A formal evidence channel taxonomy with jurisdiction sets and
    trust ceilings, implemented in \texttt{jugeo.evidence.channels} and
    \texttt{jugeo.orchestration.mixed\_evidence\_routing}.
  \item A splitting-and-dispatch pipeline that decomposes mixed VCs into
    pure-theory fragments and dispatches them in parallel to specialised
    backends.
  \item A reassembly algebra based on the trust-meet operation of the
    \JuGeo{} trust algebra, ensuring that the assembled bundle is
    conservative in its trust claim.
  \item A formal \emph{Routing Soundness Theorem} (Theorem~\ref{thm:routing-soundness})
    and its \LEAN{} formalisation, guaranteeing that fragment-level
    validity implies condition-level validity.
  \item An experimental evaluation showing \routeacc{} routing accuracy
    and $8\times$ latency improvement over monolithic dispatch on a
    \numroute{}-program benchmark.
\end{enumerate}

\paragraph{Future work.}
Several directions remain open.  First, the fragment classifier does
not yet handle \emph{quantified} mixed conditions; extending it to
AUFLIA and EUF would cover more industrial VCs.  Second, the load
balancing strategy (\texttt{LOAD\_BALANCED}) currently uses a static
semantic-weight function; an adaptive variant that updates weights
based on observed backend performance would improve throughput in
long-running verification campaigns.  Third, the trust-tier reassembly
currently applies only the \texttt{TRUST\_MEET} strategy for soundness;
a formal characterisation of conditions under which \texttt{TRUST\_JOIN}
is also sound would expand the set of verification conditions that can
receive higher trust tiers without renegotiation.

\paragraph{Integration.}
Mixed evidence routing is implemented as a drop-in component of the
\JuGeo{} orchestration layer.  Existing users of
\texttt{jugeo.orchestration} gain theory-aware dispatch automatically
by enabling \texttt{mixed\_routing=True} in the orchestrator
configuration; the router is backward-compatible with monolithic VCs,
which are treated as single-fragment conditions and routed directly.

% ── Bibliography ──────────────────────────────────────────────────
\begin{thebibliography}{9}

\bibitem{Nelson1979}
G.~Nelson and D.~C.~Oppen.
\newblock Simplification by cooperating decision procedures.
\newblock \textit{ACM Trans.\ Program.\ Lang.\ Syst.}, 1(2):245--257, 1979.

\bibitem{Martinez2019}
G.~Mart\'{\i}nez, D.~Ahman, V.~Dumitrescu, N.~Giannarakis, C.~Hawblitzel,
C.~Hrit\c{c}u, M.~Narasimhamurthy, Z.~Paraskevopoulou, C.~Pit-Claudel,
J.~Protzenko, T.~Ramananandro, A.~Rastogi, and N.~Swamy.
\newblock Meta-{F}$^\star$: Proof automation with SMT, tactics, and
metaprograms.
\newblock \textit{J.\ Funct.\ Program.}, 29:e8, 2019.

\bibitem{dMoura2008}
L.\ de~Moura and N.~Bj\o{}rner.
\newblock {Z3}: An efficient {SMT} solver.
\newblock In \textit{Proc.\ TACAS}, pages 337--340. Springer, 2008.

\bibitem{Bjorner2015}
N.~Bj\o{}rner, A.~Gurfinkel, K.~McMillan, and A.~Rybalchenko.
\newblock {Horn} clause solvers for program verification.
\newblock In \textit{Fields of Logic and Computation II}, pages 24--51.
Springer, 2015.

\bibitem{JuGeo01}
H.~Young.
\newblock Semantic sites and judgment geometry.
\newblock \textit{Judgment Geometry Series}, Paper~01, 2024.

\bibitem{JuGeo04}
H.~Young.
\newblock Trust algebra for judgment geometry.
\newblock \textit{Judgment Geometry Series}, Paper~04, 2024.

\bibitem{JuGeo05}
H.~Young.
\newblock Fragment-aware {VC} routing.
\newblock \textit{Judgment Geometry Series}, Paper~05, 2024.

\bibitem{JuGeo22}
H.~Young.
\newblock Treaty memory: Persistent multi-module contract negotiation.
\newblock \textit{Judgment Geometry Series}, Paper~22, 2025.

\end{thebibliography}

\end{document}
